\documentclass[12pt]{article}
\usepackage{latexblog}

% ============================================================
% Blog Metadata
% ============================================================
\blogtitle{搭建Nexus私服}
\blogdate{2017-06-25}
\blogtags{Docker, 云计算, 容器化}
\bloglang{zh}
\blogsummary{}

\begin{document}
\makeblogtitle

我们安装\href{https://store.docker.com/community/images/sonatype/nexus3}{Nexus}来作为我们的Docker镜像仓库。

\begin{Shaded}
\begin{Highlighting}[]
\ExtensionTok{sodu}\NormalTok{ docker pull sonatype/nexus3}
\FunctionTok{mkdir}\NormalTok{ /home/docker/nexus}
\FunctionTok{chown} \AttributeTok{{-}R}\NormalTok{ 200 /home/docker/nexus}
\FunctionTok{sudo}\NormalTok{ docker run }\AttributeTok{{-}d} \DataTypeTok{\textbackslash{}}
    \AttributeTok{{-}{-}name}\NormalTok{ nexus }\DataTypeTok{\textbackslash{}}
    \AttributeTok{{-}v}\NormalTok{ /home/docker/nexus }\DataTypeTok{\textbackslash{}}
    \AttributeTok{{-}p}\NormalTok{ 8081:8081 }\DataTypeTok{\textbackslash{}}
\NormalTok{    sonatype/nexus3}
\end{Highlighting}
\end{Shaded}

安装完成后，可以访问192.168.56.101:8081，登录进去后，添加一个Docker Hosted源，比如http://192.168.56.101:8081/repository/cloud-images/

我们可以把项目中需要用到的文件传到Nexus上。我们新建一个\texttt{raw}格式的repository

\begin{Shaded}
\begin{Highlighting}[]
\ExtensionTok{curl} \AttributeTok{{-}{-}fail} \AttributeTok{{-}u}\NormalTok{ admin:admin123 }\AttributeTok{{-}{-}upload{-}file}\NormalTok{ gradle{-}4.0{-}bin.zip }\StringTok{\textquotesingle{}http://192.168.56.101:8081/repository/files/\textquotesingle{}}
\end{Highlighting}
\end{Shaded}

这样就可以通过http://192.168.56.101:8081/repository/files/gradle-4.0-bin.zip 来访问我们这个文件了。


\end{document}
